\subsection{Title: Research Motivation for ESG ABSA in Indonesian Disclosures: Why Language, Regulation, and Depth Matter Now}

\subsubsection{Introduction}
Environmental, Social, and Governance (ESG) disclosure analysis has ascended from a peripheral corporate reporting concern to a central input for investment decisions, regulatory compliance, and stakeholder accountability. In the specific domain of aspect-based sentiment analysis (ABSA) applied to ESG disclosures (ESG-ABSA), the motivation rests on four interlocking pressures: rising regulatory scrutiny, growing demand for ESG transparency from investors, linguistic and cross-language challenges in bilingual Indonesian/English reports, and a persistent gap between disclosure volume and analytical depth. Synthesizing evidence across multiple streams shows that the strategic value of ESG disclosures depends not only on what is disclosed but on how it is analyzed, interpreted, and governed. This section builds a cohesive argument for why this research is timely, technically rigorous, and policy-relevant, culminating in a transition to the research background.

\subsubsection{Regulatory pressure on sustainability reporting globally and in Indonesia}
Global regulatory tightening elevates the baseline for ESG disclosures. Regulators worldwide have moved toward mandatory or enhanced disclosure frameworks, fostering comparability and accountability. For example, cross-jurisdictional trends highlight the increasing prominence of standardized ESG reporting to support market integrity and investor protection 
\cite{10.22617/tcs210472-2,10.2139/ssrn.4359928,10.1108/mf-01-2023-0020,10.1007/s40804-021-00233-z}. The European emphasis on SFDR, ESRS, and related taxonomy underscores a shift from voluntary to regulated ESG transparency, prompting firms to align disclosure content with evolving criteria and to demonstrate governance of climate-related risks \cite{10.3390/su152416860,10.1108/ara-07-2023-0201,10.22617/tcs210472-2}

Indonesia-specific regulatory developments are contemporaneous with regional convergence toward formalized ESG reporting. The Indonesian Financial Services Authority (OJK) and the Indonesia Stock Exchange (IDX) impose disclosures and reporting obligations that shape corporate communication with capital markets. OJK regulations intersect with IDX requirements, including electronic submission and public accessibility of disclosure documents via IDXNet, signaling a formal regulatory channel for ESG information and governance expectations \cite{10.22617/tcs210472-2} The Indonesian regulatory regime thus creates a structured environment in which ESG disclosures must be produced, structured, and traceable, elevating the need for robust analytical tools to interpret disclosures across filings, annual reports, and supplementary narratives \cite{10.22617/tcs210472-2}

The gap between regulatory demand and analytic execution motivates ABSA-in-ESG work. As regimes require more public, standardized ESG information, the challenge becomes not only data collection but also extracting decision-useful signals (risk, opportunity, governance strength) from unstructured text. Systematic text analytics—particularly ABSA—offers a mechanism to translate regulatory language into granular, action-oriented insights (e.g., identifying which environmental or governance aspects drive perceived risk or opportunity) \cite{10.1002/csr.70089,10.3390/su15054134,10.1108/mf-01-2023-0020}

\subsubsection{Rising investor demand for ESG transparency}
Investors increasingly rely on ESG disclosures to assess risk, resilience, and long-term value. The literature shows a robust link between high ESG transparency and investor decision-making, risk management, and performance analysis, with AI/ML methods used to enhance interpretability and predictive value of ESG signals \cite{10.69593/ajbais.v4i04.130,10.2139/ssrn.3738629,10.1108/mf-01-2023-0020,10.22617/tcs210472-2}. ABSA enables investors to move beyond aggregate ESG scores to consider sentiment and signals at the level of specific ESG aspects (e.g., emissions, governance quality, labor practices), which aligns with institutional due diligence and risk monitoring practices \cite{10.1002/csr.70089,10.69593/ajbais.v4i04.130,10.1108/mf-01-2023-0020}

The role of ML and NLP in ESG analytics is increasingly emphasized in governance and investment research. Systematic reviews and surveys indicate that ML/NLP are instrumental in risk detection, controversy monitoring, and trend analysis within ESG disclosures, often accompanied by explainable AI (XAI) techniques to satisfy stakeholder demands for transparency and accountability \cite{10.1002/csr.70089,10.2139/ssrn.3738629,10.1108/mf-01-2023-0020}. For Indonesian and regional markets, ML-enhanced ESG analysis can improve the reliability and timeliness of disclosures used by analysts and investors, particularly when data volumes are large and heterogeneous \cite{10.69593/ajbais.v4i04.130,10.1108/mf-01-2023-0020}

Cross-linguistic and cross-market investor needs amplify the demand for ABSA that can capture nuance across stakeholders. Investors require not only sentiment polarity but also the specific aspects driving sentiment, which ABSA is designed to deliver. Cross-lingual ABSA literature highlights the importance of multilingual and cross-language capability to support international investment analysis, as global portfolios assess Indonesian disclosures alongside English-language reports \cite{10.1016/j.inffus.2025.103073,10.1111/eufm.12509,10.48550/arxiv.2511.00024}. The intersection of ABSA with cross-lingual transfer learning is especially relevant for Indonesian bilingual reporting contexts.

\subsubsection{The challenge of bilingual Indonesian/English sustainability reports}
Bilingual ESG disclosures compound language complexity in automated analysis. The Indonesian market often features sustainability disclosures in both Indonesian and English, either within same documents or across companion reports. This bilingual setting raises unique NLP challenges, including code-switching, mixed-language terminology, and varying regulatory lexicons. Cross-lingual ABSA research documents successful strategies using multilingual pre-trained models and cross-language alignment to enable unified analysis across languages, illustrating that model architectures and data curation choices are critical for robust ABSA in bilingual contexts \cite{10.1016/j.inffus.2025.103073,10.1111/eufm.12509,10.48550/arxiv.2511.00024}

Cross-lingual ABSA advances provide a blueprint for Indonesian ESG text analytics. Surveys of cross-lingual ABSA approaches show that combining translation, alignment-based labeling, and multilingual pre-trained language models (e.g., mBERT, XLM-R) can deliver effective cross-language sentiment and aspect extraction, particularly when annotated multilingual corpora exist or when zero-shot transfer is employed \cite{10.1016/j.inffus.2025.103073}. In environments where English-language ESG disclosures are prevalent alongside local Indonesian content, such cross-lingual capabilities become a core technical requirement for ABSA pipelines \cite{10.1016/j.inffus.2025.103073,10.1111/eufm.12509}

The practicalities of bilingual ESG reporting—readability, regulatory alignment, and content coverage—also feed into the ABSA motivation. Studies of CSR/ESG reporting using NLP demonstrate that readability and content quality affect perceived credibility and the likelihood of greenwashing detection, underscoring the need for ABSA systems that can dissect environmental, social, and governance aspects across languages with interpretability \cite{10.1111/eufm.12509}. Moreover, regulatory mandates requiring standardized disclosures heighten the need for ABSA to support regulatory reporting and compliance-monitoring tasks by indexing topics and sentiments to specific regulatory themes \cite{10.1007/s40804-021-00233-z,10.2139/ssrn.3738629}

\subsubsection{The gap between disclosure volume and analytical depth}
Escalating disclosure volumes outpace analytical depth. As ESG reporting has grown in scope and frequency, the sheer quantity of textual disclosures—annual sustainability reports, regulatory fillings, press statements—has expanded dramatically. However, traditional analyses rely on coarse-grained scores or manual content coding, which are time-consuming, prone to inconsistency, and ill-suited for scalable investment diligence. A wave of research demonstrates that while disclosure volume has risen, there is a persistent gap in depth: the ability to extract granular, evidence-based insights at the aspect level, detect controversy or greenwashing, and translate textual signals into decision-relevant metrics \cite{10.1002/csr.70089,10.1111/eufm.12509,10.2139/ssrn.3738629,10.1108/mf-01-2023-0020}

AI-driven ABSA and related NLP methods address this depth gap by enabling fine-grained extraction of target aspects, sentiment polarity, and category mappings from large corpora. ABSA supports more nuanced risk and opportunity assessment by tying sentiments to specific ESG topics (e.g., emissions intensity, supply chain governance, diversity policies). Systematic mappings across studies show that ML/NLP are increasingly used to predict ESG performance signals, identify controversies, and assist in reporting alignment with standards, while XAI methods help stakeholders understand model decisions—mitigating concerns about opacity and bias \cite{10.1002/csr.70089,10.2139/ssrn.3738629,10.1108/mf-01-2023-0020}

The Indonesian context compounds these depth concerns with regulatory and market-specific factors. In Indonesia, the combination of OJK and IDX disclosures creates a data ecosystem where standardization is evolving, and where language diversity adds to the complexity of extracting unified signals. Bridging the depth gap requires ABSA pipelines that can operate on multilingual ESG narratives, handle regulatory-specific lexicons, and produce interpretable outputs that regulators and investors can trust. Evidence from related reviews and empirical investigations supports the utility of ML/NLP to map and quantify ESG content, risk signals, and governance themes in corporate disclosures, while pinpointing challenges such as data standardization, model explainability, and potential bias \cite{10.69593/ajbais.v4i04.130,10.2139/ssrn.3738629,10.1108/mf-01-2023-0020}

\subsubsection{Synthesis: why ABSA is well-suited to ESG disclosure analysis in Indonesia}
ABSA aligns with the need for granular, aspect-level insights. By simultaneously identifying target terms, aspects, opinions, and sentiment polarity, ABSA yields a structured representation of ESG narratives that reflect stakeholder concerns, regulatory focus areas, and governance signals. This granularity is essential for both investors and regulators who require precise indicators of environmental risks, social commitments, and governance practices, particularly where narratives contain nuanced or conflicting cues. Empirical works in ESG/ML contexts demonstrate ABSA’s capacity to reveal detail beyond coarse sentiment or document-level scores, including cross-aspect interactions and topic-level patterns \cite{10.3390/math12193119,10.1002/csr.70089,10.1111/eufm.12509,10.1109/access.2023.3348342,10.2139/ssrn.3738629}

Multilingual and cross-lingual capabilities are crucial in bilingual Indonesian/English reports. Cross-lingual ABSA research illustrates practical methods to extend analysis across languages using multilingual pre-trained models, translation strategies, and alignment techniques, enabling cohesive benchmarking and comparability across Indonesian disclosures and international counterparts \cite{10.1016/j.inffus.2025.103073,10.1111/eufm.12509,10.48550/arxiv.2511.00024}. This is particularly relevant for Indonesian market participants who must assess local disclosures in tandem with global ESG narratives.

Explainability and regulatory relevance. The integration of XAI methods (e.g., SHAP, LIME) with ABSA outputs addresses demands for transparency in ESG scoring and governance decisions. Regulators, investors, and auditors favor interpretable models that connect textual signals to concrete ESG dimensions and regulatory criteria. Systematic reviews of ESG analytics emphasize the value of explainability, bias mitigation, and robust evaluation for trustworthiness in financial contexts \cite{10.1002/csr.70089,10.2139/ssrn.3738629,10.1108/mf-01-2023-0020}

Regulatory and investor alignment with depth-driven analytics. The convergence of global regulatory frameworks and investor expectations creates a policy-relevant research space in which ABSA-based ESG analytics can support both compliance monitoring and performance assessment. Existing literature documents regulatory-driven disclosure obligations, the transformation of data into decision-useful signals, and the increasing role of ML/NLP in ESG governance and investment decision-making \cite{10.1007/s40804-021-00233-z,10.69593/ajbais.v4i04.130,10.2139/ssrn.3738629,10.1108/mf-01-2023-0020}

\subsubsection{Concluding transition to research background}
The foregoing synthesis establishes that ESG disclosure analysis is currently propelled by three interdependent forces: (a) escalating regulatory expectations and formalization of ESG reporting in Indonesia and globally; (b) rising demand from investors for detailed, actionable ESG signals; and (c) the technical necessity to handle bilingual Indonesian/English disclosures with scalable, interpretable ABSA techniques that close 

\subsubsection{The gap between disclosure volume and analytical depth}. These dynamics jointly justify the development of ABSA-enabled ESG analysis tailored to Indonesian disclosures, with attention to multilingual capabilities, regulatory lexicons, and explainable outputs. The transition to the research background now follows: we outline the existing empirical and methodological landscape, identify gaps specific to ABSA in Indonesian ESG reporting, and articulate the research questions, data sources, and methodological approach that will be pursued.

% References (IEEE style)

% Khak et al. (2025) Khak et al., Bridging Finance and AI: A Comprehensive Survey of Large Language Models in Financial Applications, 2025
% (Heever et al., 2024). Heever et al., Understanding Public Opinion towards ESG and Green Finance with the Use of Explainable Artificial Intelligence, Math., 2024
% (Xu, 2022). Xu, AI Fairness in the Financial Industry: A Machine Learning Pipeline Approach, 2022
% (Adeoye et al., 2024). Adeoye et al., Artificial Intelligence in ESG investing: Enhancing portfolio management and performance, 2024
% (Šmíd & Král, 2025). Šmíd and Král, Cross-lingual aspect-based sentiment analysis: A survey on tasks, approaches, and challenges, Information Fusion, 2025
% (Seow, 2025). Seow, Transforming ESG Analytics With Machine Learning: A Systematic Literature Review Using TCCM Framework, 2025
% (Miglionico, 2022). Miglionico, The Use of Technology in Corporate Management and Reporting of Climate-Related Risks, European Business Organization Law Review, 2022
% (Gupta et al., 2024). Gupta et al., A Natural Language Processing Model on BERT and YAKE Technique for Keyword Extraction on Sustainability Reports, IEEE Access, 2024
% (Gorovaia & Makrominas, 2024). Gorovaia and Makrominas, Identifying greenwashing in corporate-social responsibility reports using natural-language processing, European Financial Management, 2024
% (Bogdan et al., 2023). Bogdan et al., A Streamline Sustainable Business Performance Reporting Model by an Integrated FinESG Approach, Sustainability, 2023
% (Lahmar et al., 2023). Lahmar et al., ESG in the Financial Industry: What Matters for Rating Analysts?, SSRN, 2023
% (Hang, 2025). Hang, Chitchat with AI: Understand the supply chain carbon disclosure of companies worldwide through Large Language Model, arXiv, 2025
% (Alshaikh et al., 2024). Alshaikh et al., BERT-Based Model for ABSA for Arabic Open-Ended Survey Responses: A Case Study, IEEE Access, 2024
% (Juthi et al., 2024). Juthi et al., Sustainable Finance and Data Analytics: A Systematic Review of ESG Data in Investment Decisions, Ajbais, 2024
% (Zhang & Zhang, 2023). Zhang and Zhang, Renovation in ESG research: ML applications, Asian Review of Accounting, 2023
% (Moreno & Caminero, 2020). Sultana and Zeya, Advancing Firm-Specific ESG Sentiment Analysis: A Machine Learning Approach to Clustering, Prediction, and Financial Performance Implications, Business Strategy & Development, 2025
% (Sultana & Zeya, 2025). Rutherford, Aspect-Level Obfuscated Sentiment in Thai Financial Disclosures and Its Impact on Abnormal Returns, arXiv, 2025
% (Raghupathi et al., 2023). Raghupathi et al., Understanding Corporate Sustainability Disclosures from the SEC Filings, Sustainability, 2023
% (Rutherford, 2025). 10.48550/arxiv.2511.13481, Rutherford, 2025
% (McBride et al., 2022). 10.1108/mf-12-2021-0586, McBride et al., 2022
% (Rane et al., 2024). 10.70593/978-81-981367-4-9_3, Rane et al., 2024
% (Ding, 2024). 10.33423/jabe.v26i4.7181, Ding, 2024
% (И.Р., 2024). 10.53555/kuey.v30i5.4019, И.Р., 2024
% (Srivastava & Anand, 2023). 10.1108/mf-01-2023-0020, Srivastava and Anand, 2023
% ("The Bond Market in Indonesia:", 2021). 10.22617/tcs210472-2, Rane et al., 2024.

% Annotated bibliography (machine-readable XML)